\subsection{R discovery when borrower risk and selection are hidden}
\label{exec13:latent-r}

The preceding PC/IDA boundary assumes a causally sufficient system. That
assumption is consequential in a servicing application: private borrower
distress can affect both the bank's contact decision and subsequent recovery,
while inclusion in an observed account panel can depend on both. A larger
regression or a more elaborate graph display does not by itself remove either
problem. The R comparison must therefore ask what information survives latent
common causes and selection, which object an algorithm returns, and which
additional argument would turn that object into a decision.

The selected implementation is again \texttt{pcalg} 2.7-12 at the preserved
September 2024 CRAN mirror revision. FCI, RFCI, FCI+, conservative and majority
variants, and the Joint Causal Inference options are related operations in
this package. They are not six independent software implementations. The
bounded R search also returns methodological papers, a Python software paper
and other-language leads; their presence in a search result does not establish
an R implementation. The following claims concern inspected pinned manuals,
implementation paths and authentic method papers. They do not constitute a
complete census of every R extension or an executed package comparison
\citep{exec13_latentr_primary}.

\paragraph{What a latent-variable graph says.}
Let $X$ denote the recorded variables, $L$ unrecorded variables, and $S$ the
variables on which inclusion is conditioned. For an underlying acyclic causal
graph, a maximal ancestral graph can represent the conditional-independence
structure left after marginalizing $L$ and conditioning on $S$. A partial
ancestral graph, or PAG, records features shared by compatible ancestral
graphs. The marks carry information: an arrowhead at a vertex rules out its
ancestry of the opposite vertex and the selection variables under the stated
model; a tail carries an ancestral restriction; a circle leaves an endpoint
unresolved. In particular, selection can make the interpretation of a tail
refer to ancestry of selection, rather than a uniquely established direct
causal effect on the other recorded variable. An estimated PAG is neither a
single fully specified structural model nor a numerical intervention estimate
\citep{exec13_latentr_primary,exec13_latentr_colombo2012}.

The distinction is stronger for RFCI. In the definition of an FCI-PAG, an
adjacency asserts dependence for all eligible conditioning subsets of the
recorded variables, conditional on selection. In an RFCI-PAG the corresponding
requirement concerns subsets of the adjacent sets. This weaker statement can
leave additional edges and can represent more than one Markov equivalence
class. The causal endpoint information is sound in the oracle theorem under
faithfulness, but it is incorrect to relabel the RFCI output as the complete
FCI equivalence class. A common serialization class named \texttt{fciAlgo}
does not erase this difference. The proposed artifact must retain the
algorithm, version and object semantics as well as the matrix
\citep{exec13_latentr_colombo2012,exec13_latentr_primary}.

\paragraph{A bank decision that the observed table cannot settle.}
Consider a binary contact action $A$ and binary recovery $Y$. Three authored
mechanisms produce exactly the same strictly positive observed table:
\[
\begin{array}{c|cc}
 &Y=0&Y=1\\\hline
 A=0&2/5&1/10\\
 A=1&1/10&2/5
\end{array}
\]
In the first mechanism $A$ is randomized with probability $1/2$, and
$P(Y=1\mid\operatorname{do}(A=a))=1/5+3a/5$. Its effect is $3/5$.
In the second, hidden $U$ is a fair binary variable. Conditional on $U$,
action and recovery are independent, with
\[
 P(A=1\mid U=u)=1/10+4u/5,\qquad
 P(Y=1\mid U=u)=1/8+3u/4.
\]
All eight $(U,A,Y)$ cells have positive probability. Marginalization gives
the same displayed table, but setting $A$ does not alter $U$ or $Y$: both
intervention means are $1/2$ and the effect is zero. Thus an observed
conditional contrast of $3/5$ is compatible with no contact effect
\citep{exec13_latentr_exact}.

For the third mechanism, $A$ and $Y$ are independent in the population,
with $P(A=1)=2/5$ and $P(Y=1)=1/3$. Inclusion probabilities for
$(A,Y)=(0,0),(0,1),(1,0),(1,1)$ are respectively
$1/5,1/10,3/40,3/5$. Multiplying by the population cell probabilities
and normalizing by $P(S=1)=1/5$ again gives the displayed table.
The population effect is zero; the selected association
does not identify the unselected population effect. The selection graph has
$A\to S\leftarrow Y$. It has a different interpretation from the hidden
common-cause graph $A\leftarrow U\to Y$, even though this observed table
cannot distinguish them from the direct-effect mechanism
\citep{exec13_latentr_exact}.

If a recovery is assigned a value of USD~100 and a contact costs USD~20, the
three stipulated population net values are USD~40, USD~$-20$ and USD~$-20$.
These are illustrative prices and mechanisms, not bank measurements. The
three values demonstrate unresolved decision consequences; they are not a
sharp identified set over all possible causal models. A useful product keeps
the conflicting explanations and asks which restriction, assignment evidence
or intervention could distinguish them. Selecting the graph whose implied
profit is most attractive cannot supply that missing information
\citep{exec13_latentr_exact}.

\paragraph{How the algorithms spend information and computation.}
FCI begins with an adjacency search and then considers separating sets that
can lie outside a node's immediate neighborhood. Its Possible-D-SEP stage
can remove edges the initial search retained, after which orientation rules
refine endpoint marks. RFCI omits that global stage and instead performs
additional local tests for unshielded triples and discriminating paths.
This is a trade between computational work and informativeness, not a claim
that RFCI finds exactly the same graph more cheaply on every input. The
oracle soundness result requires a faithful underlying DAG and correct
conditional-independence information. Neither a software return code nor a
large nominal sample size establishes these premises
\citep{exec13_latentr_colombo2012,exec13_latentr_primary}.

The high-dimensional consistency result is also specific. It assumes IID
copies of the selected observed law, multivariate Gaussianity, faithfulness,
restrictions on graph-neighborhood growth, and nonzero partial correlations
bounded away from a shrinking detection boundary while remaining bounded
away from one. The FCI result uses a stronger restriction on the growth of
Possible-D-SEP sets. Its conclusion uses a sequence of tuning levels; it is
not a finite-sample graph-confidence guarantee at the user's chosen
\texttt{alpha}. Repeated monthly observations on the same borrowers do not
become IID merely because they are supplied to the same function
\citep{exec13_latentr_colombo2012}.

The per-test level should consequently travel as a test parameter, not a
confidence attached to the entire graph. For an explicitly hypothetical
family of 100 independent true-null tests at level $.05$, the probability
of at least one rejection is $1-.95^{100}$. That elementary calculation
does not describe FCI's adaptive testing error, but it exposes the unsupported
step from a per-test parameter to a simultaneous claim. Graph uncertainty,
orientation ambiguity, model restrictions and sampling uncertainty need
separate treatment before an effect set can receive a coverage label
\citep{exec13_latentr_exact,exec13_latentr_colombo2012}.

FCI+ supplies a different search strategy. Its method paper bounds the
number of oracle independence tests by order $N^{2(k+2)}$ for $N$ observed
variables and a fixed bound $k$ on node degree. The result presumes the
faithful acyclic model and the stated oracle conditions while allowing
latent variables and selection. Polynomial order at fixed $k$ is not a
practical runtime promise, and the paper notes that sample versions can
use larger conditioning sets and lose power. The proposal should compare
FCI+ as an alternative under the same generated regimes, without translating
the complexity result into a bank-scale performance claim
\citep{exec13_latentr_claassen2013}.

\paragraph{Options are scientific choices, not cosmetic settings.}
The pinned implementation exposes normal, anytime and adaptive-anytime
FCI. A finite \texttt{m.max} restricts conditioning-set size. A distinct
parameter, \texttt{pdsep.max}, is checked against the size of a candidate
Possible-D-SEP set: when that size exceeds the limit, the selected code skips
the corresponding search. It does not simply enumerate all subsets up to
that order. A job that reaches the limit therefore needs an explicit
restricted-search disposition. A faster answer is not automatically
equivalent to an unrestricted run
\citep{exec13_latentr_primary}.

Missing conditional-independence results require equal care. In the inspected
Possible-D-SEP path, \texttt{NAdelete=TRUE} maps a missing test result to
one and can remove an edge; false maps it to zero and retains the edge in
that path. This is a software rule, not evidence that the missing result
establishes independence or dependence. The proposed adapter should retain
the missing-test count and reason and apply an explicit admission policy.
It should not silently use either behavior to manufacture certainty.
Hard-coded gaps and edges likewise encode supplied restrictions; their
provenance and sensitivity belong in the scientific result
\citep{exec13_latentr_primary}.

Conservative or majority-rule orientation reduces particular ambiguities
but does not imply full order invariance. The order-independent methods
paper distinguishes stable skeletons, decisions about unshielded triples
and the subsequent orientation rules. RFCI can remove further edges using
order-dependent separating sets, so even a stable initial skeleton does not
make its final skeleton fully invariant. The package manual retains this
warning. Permuting the input columns and translating back to the original
labels is therefore a useful proposed diagnostic, not a theorem that the
selected route must already pass
\citep{exec13_latentr_colombo2014,exec13_latentr_primary}.

Selection and context options must be admitted together. FCI and FCI+
default to allowing selection bias. Setting \texttt{selectionBias=FALSE}
imposes its absence and disables the selected selection-related orientation
rules. In this pinned implementation nondefault JCI assumptions require the
no-selection setting. A branch or policy-regime label is not automatically
an exogenous, unconfounded randomized context. The proposed workflow must
ask why that context restriction is justified and whether account inclusion
is compatible with it. The wider cyclic and JCI theoretical claims require
their own source and method review; this comparison does not inherit them
from an exposed argument name
\citep{exec13_latentr_primary}.

\paragraph{Two concrete boundary defects change the reuse decision.}
For \texttt{amat.pag}, the matrix entry refers to the endpoint at the
column variable: one is a circle, two an arrowhead and three a tail. This
differs from the package's DAG/CPDAG convention. The pinned
\texttt{pcalg2dagitty} implementation has explicit branches for seven
nonzero endpoint pairs but omits the two tail--circle cases. In particular,
the documented input
\[
 \begin{pmatrix}0&1\\3&0\end{pmatrix}
\]
represents a tail at the first vertex and a circle at the second. A static
trace through this two-node conversion adds neither an edge nor an isolated
node and constructs \texttt{pag\{\}} before the parser call. This is a
source-level counterexample to unconditional lossless conversion, not a
recorded native parser response. The proposed route must explicitly support
both cases or refuse them before conversion, then verify labels, all endpoints
and round-trip identity in an isolated runtime qualification
\citep{exec13_latentr_primary,exec13_latentr_exact}.

FCI+ has a different provenance problem. Its manual describes the general
result object as containing separating sets, but the inspected return
constructor places empty lists or vectors in several diagnostic fields,
including separating sets and test counts, with a placeholder
\texttt{pMax}. Empty fields must not be reported as a measured zero-test
search or as a complete explanation of why an edge disappeared. RFCI's
reported initial-skeleton test counts also do not enumerate all of its
later local tests. A scientifically auditable route needs instrumentation
of the conditional-independence interface or an explicit diagnostic-coverage
limitation. This work has inspected those code paths; it has not installed
that instrumentation
\citep{exec13_latentr_primary}.

A representative-graph conversion cannot repair either scientific uncertainty
or missing provenance. The inspected \texttt{pag2magAM} documentation
constructs one compatible MAG, subject to a restriction on new edges into
the requested node and the absence of selection variables. Its output is a
representative used within a specified graphical criterion, not evidence
that the selected graph is the true bank mechanism. In particular, the
existence of one convenient adjustment or effect calculation does not settle
all models represented by the original discovery result
\citep{exec13_latentr_primary}.

\paragraph{A restricted integration path and its alternatives.}
The preferred addition to the proposal is an isolated R FCI/RFCI comparison
route that preserves PAG versus RFCI-PAG semantics and all endpoint marks.
FCI is the richer reference; RFCI is the explicitly less informative
alternative when the proposed workload makes the full search unattractive.
FCI+ remains a candidate subject to complete diagnostic accounting and
matched resource evaluation. The initial PC/IDA path remains appropriate
only when its separate causal-sufficiency assumptions are defended. The
existing Python and Java comparisons offer independent implementation
candidates, but equal algorithm names do not establish equal defaults,
conversion conventions or statistical outputs.

Direct adoption of an arbitrary returned graph as a bank decision is rejected.
A wholesale port is also not yet justified: it would inherit a large
scientific verification burden, and moving the same logic into another
language would not resolve latent information or selection. The immediate
reuse choice is a narrow adapter with explicit test, search and diagnostic
contracts. The package's previously inspected GPL declaration and dependency
inventory remain relevant to any distribution decision; this scientific
selection grants neither a new licensing conclusion nor deployment approval.

The proposed qualification should begin with independently specified
independence oracles and small known graphs, then use generated samples with
hidden causes, selection, weak associations and missing tests. It must retain
all requests, budget exits and graph ambiguities; compare FCI, RFCI and FCI+
on the same declared laws; and separate native numerical behavior from
serialization correctness. Endpoint pairs, isolated nodes, invalid graph
types, column permutations, diagnostic completeness and search truncation
need individual fixtures. Only after those results should a decision study
evaluate whether a justified effect or refusal changes contact policy value.
No package execution, finite-sample calibration or accepted R adapter is
claimed by the present source review.
\label{exec13:latent-r-end}
